
\sgaStableTarget{sga6.e10.sec.7}\section*{7. Appendix: Specialization in Intersection Theory}
\begin{center}
\underline{by A. Grothendieck}
\end{center}

\sgaStableTarget{sga6.e10.lem.7-1}\textbf{Lemma 7.1.} \uline{Let} @@M0@@ \uline{be a regular closed immersion} (VII 1.4)\uline{, of bounded codimension, so that} @@M1@@ \uline{is a perfect morphism} (VII 1.9)\uline{, and one can therefore define the morphism} (IV 2.12)
@@M2@@
\uline{Let, on the other hand,} @@M3@@ \uline{be the conormal sheaf of} @@M4@@\uline{, which is a locally free sheaf of finite type on} @@M5@@ (VII 1.6)\uline{, so that one can form}
@@M6@@
\uline{Consider also the direct-image homomorphism} (IV 2.11)
@@M7@@
\uline{Then one has the formula}
@@M8@@
\uline{where the product on the right is taken with respect to the} @@M9@@\uline{-module structure on} @@M10@@ \uline{defined in} (IV 2.10).

The proof is, essentially without change, the same as that of the same formula, written for @@M11@@, @@M12@@, established in (VII 2.7).

\sgaStableTarget{sga6.e10.cor.7-2}\textbf{Corollary 7.2.} \uline{Let} @@M13@@ \uline{be a flat morphism, with} @@M14@@ \uline{a trait, i.e., the spectrum of a discrete valuation ring} @@M15@@, @@M16@@ \uline{the closed point of} @@M17@@, @@M18@@ \uline{its generic point, and} @@M19@@ \uline{the canonical inclusion. Then} @@M20@@ \uline{has finite Tor-dimension, so} (IV 2.12) \uline{one can define}
@@M21@@
